% GENERATED FILE. Edit canonical lessons, then run npm run generate:books.
% canonical-source-hash: fnv1a64:ef4609d081b2f8a5
% canonical-lessons: ML-C215-pantu, ML-C215-kalikkaran, ML-C215-tim, ML-C215-samanila, ML-C215-hobi

\chapter{Ball, Player, Team, Draw, Hobby}
\label{ch:ml-ball-player-team-draw-hobby}
\hlchaptermodality{215}

\begin{chapteropening}
\textbf{By the end of this chapter:} \emph{I can say a ball, a player, a team, a draw and a hobby.}

Complete the last lesson of chapter 215: I can say a ball, a player, a team, a draw and a hobby.
\end{chapteropening}

\section[\texorpdfstring{\ml{പന്ത്} (pantŭ)}{പന്ത് (pantŭ)}]{\ml{പന്ത്} (pantŭ) — a ball}

\label{lesson:ML-C215-pantu}



\begin{quote}
Before the new one: say the Malayalam for a switch, then the Malayalam for a saw.
\end{quote}

\subsection*{You'll want to know: \ml{പന്ത്}}
\textbf{\ml{പന്ത്}} — \emph{pantŭ} — \textquotedblleft{}a ball\textquotedblright{}.

\begin{cousinweb}
\textbf{\ml{പന്ത്}} (\emph{pantŭ}) is native Dravidian, one word with Tamil \emph{pantu} and Telugu \emph{banti} (DEDR 3925).
\end{cousinweb}

\begin{grammarlens}[title={What you've built}]
One more word for this chapter.
\end{grammarlens}

\subsection*{Your turn}
\begin{itemize}
\raggedright
  \item \emph{Say it:} \emph{pantŭ}
  \item \emph{Say it:} \emph{pantŭ}, once more
  \item \emph{Say it:} say \emph{pantŭ}
  \item \emph{Recall:} say the Malayalam for a switch, then the Malayalam for a saw, then say \emph{pantŭ} again
\end{itemize}

\begin{tcolorbox}[breakable,colback=teal!4,colframe=teal!35!black,title={Before you move on}]
What does \ml{പന്ത്} mean? (\textquotedblleft{}A ball\textquotedblright{}.) Say it once more.
\end{tcolorbox}
\section[\texorpdfstring{\ml{കളിക്കാരൻ} (kaḷikkāran)}{കളിക്കാരൻ (kaḷikkāran)}]{\ml{കളിക്കാരൻ} (kaḷikkāran) — a player}

\label{lesson:ML-C215-kalikkaran}



\begin{quote}
Before the new one: say the Malayalam for a saw, then the Malayalam for a ball.
\end{quote}

\subsection*{You'll want to know: \ml{കളിക്കാരൻ}}
\textbf{\ml{കളിക്കാരൻ}} — \emph{kaḷikkāran} — \textquotedblleft{}a player\textquotedblright{}.

\begin{cousinweb}
\textbf{\ml{കളിക്കാരൻ}} (\emph{kaḷikkāran}) is \emph{kaḷi}, \textquotedblleft{}play\textquotedblright{} (DEDR 1374), with \emph{-kkāran}, the ending for a person who does something: \textquotedblleft{}one who plays\textquotedblright{}. DEDR 1374 compares Sanskrit \emph{kalyā} and \emph{kaḍ}.
\end{cousinweb}

\begin{grammarlens}[title={What you've built}]
One more word for this chapter.
\end{grammarlens}

\subsection*{Your turn}
\begin{itemize}
\raggedright
  \item \emph{Say it:} \emph{kaḷikkāran}
  \item \emph{Say it:} \emph{kaḷikkāran}, once more
  \item \emph{Say it:} say \emph{kaḷikkāran}
  \item \emph{Recall:} say the Malayalam for a saw, then the Malayalam for a ball, then say \emph{kaḷikkāran} again
\end{itemize}

\begin{tcolorbox}[breakable,colback=teal!4,colframe=teal!35!black,title={Before you move on}]
What does \ml{കളിക്കാരൻ} mean? (\textquotedblleft{}A player\textquotedblright{}.) Say it once more.
\end{tcolorbox}
\section[\texorpdfstring{\ml{ടീം} (ṭīṁ)}{ടീം (ṭīṁ)}]{\ml{ടീം} (ṭīṁ) — a team}

\label{lesson:ML-C215-tim}



\begin{quote}
Before the new one: say the Malayalam for a ball, then the Malayalam for a player.
\end{quote}

\subsection*{You'll want to know: \ml{ടീം}}
\textbf{\ml{ടീം}} — \emph{ṭīṁ} — \textquotedblleft{}a team\textquotedblright{}.

\begin{cousinweb}
\textbf{\ml{ടീം}} (\emph{ṭīṁ}) $\leftarrow$ English \emph{team}.
\end{cousinweb}

\begin{grammarlens}[title={What you've built}]
One more word for this chapter.
\end{grammarlens}

\subsection*{Your turn}
\begin{itemize}
\raggedright
  \item \emph{Say it:} \emph{ṭīṁ}
  \item \emph{Say it:} \emph{ṭīṁ}, once more
  \item \emph{Say it:} say \emph{ṭīṁ}
  \item \emph{Recall:} say the Malayalam for a ball, then the Malayalam for a player, then say \emph{ṭīṁ} again
\end{itemize}

\begin{tcolorbox}[breakable,colback=teal!4,colframe=teal!35!black,title={Before you move on}]
What does \ml{ടീം} mean? (\textquotedblleft{}A team\textquotedblright{}.) Say it once more.
\end{tcolorbox}
\section[\texorpdfstring{\ml{സമനില} (samanila)}{സമനില (samanila)}]{\ml{സമനില} (samanila) — a draw, a tie}

\label{lesson:ML-C215-samanila}



\begin{quote}
Before the new one: say the Malayalam for a player, then the Malayalam for a team.
\end{quote}

\subsection*{You'll want to know: \ml{സമനില}}
\textbf{\ml{സമനില}} — \emph{samanila} — \textquotedblleft{}a draw, a tie\textquotedblright{}.

\begin{cousinweb}
\textbf{\ml{സമനില}} (\emph{samanila}) joins Sanskrit \emph{sama}, \textquotedblleft{}equal\textquotedblright{} (CDIAL 13173), to native \emph{nila}, \textquotedblleft{}standing, position\textquotedblright{} (DEDR 3675, with Tamil \emph{nil}): \textquotedblleft{}equal standing\textquotedblright{}.
\end{cousinweb}

\begin{grammarlens}[title={What you've built}]
One more word for this chapter.
\end{grammarlens}

\subsection*{Your turn}
\begin{itemize}
\raggedright
  \item \emph{Say it:} \emph{samanila}
  \item \emph{Say it:} \emph{samanila}, once more
  \item \emph{Say it:} say \emph{samanila}
  \item \emph{Recall:} say the Malayalam for a player, then the Malayalam for a team, then say \emph{samanila} again
\end{itemize}

\begin{tcolorbox}[breakable,colback=teal!4,colframe=teal!35!black,title={Before you move on}]
What does \ml{സമനില} mean? (\textquotedblleft{}A draw, a tie\textquotedblright{}.) Say it once more.
\end{tcolorbox}
\section[\texorpdfstring{\ml{ഹോബി} (hōbi)}{ഹോബി (hōbi)}]{\ml{ഹോബി} (hōbi) — a hobby}

\label{lesson:ML-C215-hobi}



\begin{quote}
Before the new one: say the Malayalam for a team, then the Malayalam for a draw.
\end{quote}

\subsection*{You'll want to know: \ml{ഹോബി}}
\textbf{\ml{ഹോബി}} — \emph{hōbi} — \textquotedblleft{}a hobby\textquotedblright{}.

\begin{cousinweb}
\textbf{\ml{ഹോബി}} (\emph{hōbi}) $\leftarrow$ English \emph{hobby}.
\end{cousinweb}

\begin{grammarlens}[title={What you've built}]
One more word for this chapter.
\end{grammarlens}

\subsection*{Your turn}
\begin{itemize}
\raggedright
  \item \emph{Say it:} \emph{hōbi}
  \item \emph{Say it:} \emph{hōbi}, once more
  \item \emph{Say it:} say \emph{hōbi}
  \item \emph{Recall:} say the Malayalam for a team, then the Malayalam for a draw, then say \emph{hōbi} again
\end{itemize}

\begin{tcolorbox}[breakable,colback=teal!4,colframe=teal!35!black,title={Before you move on}]
What does \ml{ഹോബി} mean? (\textquotedblleft{}A hobby\textquotedblright{}.) Say it once more.
\end{tcolorbox}
